\documentclass{article}
\usepackage[margin=1in]{geometry}
\usepackage{amsmath,amssymb}
\usepackage[most]{tcolorbox}
\usepackage[hidelinks]{hyperref}

\newcommand{\BookName}{Pattern Recognition and Machine Learning}
\newcommand{\BookAuthor}{Christopher M. Bishop}
\newcommand{\ChapterName}{Introduction}
\newcommand{\ExerciseID}{1.12}

\title{Chapter: \ChapterName}
\author{Ahonakpon Guy Gbaguidi}
\date{}

\begin{document}

\maketitle

\begin{center}
  \large\textbf{Exercise \ExerciseID}

  \vspace{0.5em}

  \normalsize\BookName\ - \BookAuthor
\end{center}

\begin{tcolorbox}[breakable,colback=gray!5,colframe=gray!60,title=Solution]
In this exercise, we are asked to show that

\[
\mathbb{E}[x_n x_m] = \mu^2 + I_{nm} \sigma^2
\]
where $x_1,\ldots,x_N$ are independent samples from a Gaussian distribution with mean $\mu$ and variance $\sigma^2$, and $I_{nm}$ satisfies $I_{nm} = 1$ if $n = m$ and $I_{nm} = 0$ otherwise.

For $n\neq m$, the expectation of the product is

\[
\mathbb{E}[x_n x_m] = \int_{-\infty}^{\infty} \int_{-\infty}^{\infty} x_n x_m p(x_n, x_m) \, dx_n \, dx_m
\]

We know that if $n \neq m$, then $x_n$ and $x_m$ are independent, and thus the joint probability density function can be expressed as the product of their individual densities:

\[
p(x_n, x_m) = p(x_n) p(x_m)
\]

Then, we can write:

\begin{align*}
\mathbb{E}[x_n x_m] 
&= \int_{-\infty}^{\infty} \int_{-\infty}^{\infty} x_n x_m p(x_n) p(x_m) \, dx_n \, dx_m \\
&= \int_{-\infty}^{\infty} x_n \left( \int_{-\infty}^{\infty} x_m p(x_m) \, dx_m \right) p(x_n) \, dx_n \\
&= \int_{-\infty}^{\infty} x_n \mathbb{E}[x_m] p(x_n) \, dx_n \\
&= \mathbb{E}[x_m] \int_{-\infty}^{\infty} x_n p(x_n) \, dx_n \\
&= \mathbb{E}[x_m] \mathbb{E}[x_n] \\
&= \mathbb{E}[x_n] \mathbb{E}[x_m] \\
&= \mu^2
\end{align*}

Now, if $n = m$, then we have:

\begin{align*}
\mathbb{E}[x_n x_n]
&= \mathbb{E}[x_n^2] \\
&= \text{Var}(x_n) + (\mathbb{E}[x_n])^2 \\
&= \sigma^2 + \mu^2
\end{align*}


We can conclude that:
\[
\mathbb{E}[x_n x_m] = \mu^2 + I_{nm} \sigma^2
\]

where $I_{nm} = 1$ if $n = m$ and $I_{nm} = 0$ otherwise, which completes the proof.

\paragraph{Expected maximum-likelihood mean: result (1.57).}
The estimators are
\[
\mu_{ML}=\frac{1}{N}\sum_{n=1}^{N}x_n,
\qquad
\sigma^2_{ML}=\frac{1}{N}\sum_{n=1}^{N}(x_n-\mu_{ML})^2.
\]
Both are random variables because they depend on the sample. The true
parameters $\mu$ and $\sigma^2$ are fixed, and expectations are taken over
the possible samples. For the mean,
\begin{align*}
\mathbb{E}[\mu_{ML}] &= \mathbb{E}\left[\frac{1}{N} \sum_{n=1}^{N} x_n\right] \quad \text{(definition of $\mu_{ML}$)} \\
&= \frac{1}{N} \sum_{n=1}^{N} \mathbb{E}[x_n] \quad \text{(linearity of expectation)} \\
&= \frac{1}{N} \sum_{n=1}^{N} \mu \quad \text{(each $x_n$ has mean $\mu$)} \\
&= \mu \quad \text{(there are $N$ identical terms)}
\end{align*}

\paragraph{Expected maximum-likelihood variance: result (1.58).}
Expanding the squared residual gives
\begin{align*}
\mathbb{E}[\sigma^2_{ML}]
&=\frac{1}{N}\sum_{n=1}^{N}\mathbb{E}[(x_n-\mu_{ML})^2] \\
&=\frac{1}{N}\sum_{n=1}^{N}
\left(\mathbb{E}[x_n^2]-2\mathbb{E}[x_n\mu_{ML}]
+\mathbb{E}[\mu_{ML}^2]\right).
\end{align*}
We already know that $\mathbb{E}[x_n^2]=\mu^2+\sigma^2$.
We now calculate the other two expectations explicitly.

\paragraph{1. Calculate $\mathbb{E}[x_n\mu_{ML}]$.}
Fix $n$. Since $\mu_{ML}$ includes $x_n$, they are not independent.
Substituting the sample mean and separating the term $m=n$ gives
\begin{align*}
\mathbb{E}[x_n\mu_{ML}]
&=\mathbb{E}\left[x_n\frac{1}{N}\sum_{m=1}^{N}x_m\right] \\
&=\frac{1}{N}\sum_{m=1}^{N}\mathbb{E}[x_nx_m] \\
&=\frac{1}{N}\left(\mathbb{E}[x_n^2]
+\sum_{\substack{m=1\\m\neq n}}^{N}\mathbb{E}[x_nx_m]\right) \\
&=\frac{1}{N}\left((\mu^2+\sigma^2)+(N-1)\mu^2\right) \\
&=\frac{1}{N}\left(N\mu^2+\sigma^2\right) \\
&=\mu^2+\frac{\sigma^2}{N}.
\end{align*}
There is one diagonal term and $N-1$ terms involving distinct samples.
In particular, replacing $\mathbb{E}[x_n\mu_{ML}]$ by
$\mathbb{E}[x_n]\mathbb{E}[\mu_{ML}]=\mu^2$ would be incorrect.

\paragraph{2. Calculate $\mathbb{E}[\mu_{ML}^2]$.}
The square of a sum contains all ordered pairs of indices:
\begin{align*}
\mu_{ML}^2
&=\left(\frac{1}{N}\sum_{n=1}^{N}x_n\right)
  \left(\frac{1}{N}\sum_{m=1}^{N}x_m\right) \\
&=\frac{1}{N^2}\sum_{n=1}^{N}\sum_{m=1}^{N}x_nx_m.
\end{align*}
Of the $N^2$ ordered pairs, $N$ have $n=m$ and $N(N-1)$ have $n\neq m$.
Therefore,
\begin{align*}
\mathbb{E}[\mu_{ML}^2]
&=\frac{1}{N^2}\sum_{n=1}^{N}\sum_{m=1}^{N}\mathbb{E}[x_nx_m] \\
&=\frac{1}{N^2}\left(
  \sum_{n=1}^{N}\mathbb{E}[x_n^2]
  +\sum_{n=1}^{N}\sum_{\substack{m=1\\m\neq n}}^{N}\mathbb{E}[x_nx_m]
  \right) \\
&=\frac{1}{N^2}\left(N(\mu^2+\sigma^2)+N(N-1)\mu^2\right) \\
&=\frac{1}{N^2}\left(N\mu^2+N\sigma^2+(N^2-N)\mu^2\right) \\
&=\frac{1}{N^2}\left(N^2\mu^2+N\sigma^2\right) \\
&=\mu^2+\frac{\sigma^2}{N}.
\end{align*}
Notice that $\mathbb{E}[\mu_{ML}^2]\neq(\mathbb{E}[\mu_{ML}])^2$;
their difference is $\operatorname{Var}(\mu_{ML})=\sigma^2/N$.

\paragraph{3. Substitute and simplify.}
For each $n$, the expected squared residual is
\begin{align*}
\mathbb{E}[(x_n-\mu_{ML})^2]
&=\mathbb{E}[x_n^2]-2\mathbb{E}[x_n\mu_{ML}]
  +\mathbb{E}[\mu_{ML}^2] \\
&=(\mu^2+\sigma^2)
  -2\left(\mu^2+\frac{\sigma^2}{N}\right)
  +\left(\mu^2+\frac{\sigma^2}{N}\right) \\
&=\mu^2+\sigma^2-2\mu^2-\frac{2\sigma^2}{N}
  +\mu^2+\frac{\sigma^2}{N} \\
&=(\mu^2-2\mu^2+\mu^2)
  +\sigma^2-\frac{2\sigma^2}{N}+\frac{\sigma^2}{N} \\
&=\sigma^2-\frac{\sigma^2}{N} \\
&=\frac{N-1}{N}\sigma^2.
\end{align*}
Finally, averaging these $N$ identical expectations gives
\begin{align*}
\mathbb{E}[\sigma^2_{ML}]
&=\frac{1}{N}\sum_{n=1}^{N}\mathbb{E}[(x_n-\mu_{ML})^2] \\
&=\frac{1}{N}\sum_{n=1}^{N}\frac{N-1}{N}\sigma^2 \\
&=\frac{1}{N}\cdot N\cdot\frac{N-1}{N}\sigma^2 \\
&=\boxed{\frac{N-1}{N}\sigma^2}.
\end{align*}
\end{tcolorbox}

\end{document}
